
\providecommand{\Mp}{\mathbf{Mp}}
\providecommand{\Spin}{\mathbf{Spin}}
\providecommand{\Pf}{\operatorname{Pf}}
\providecommand{\Un}{\mathbf{U}}

\chapter{The singular theta lift}
\label{ch:singular-theta-lift}

\section{The singular theta lift of \texorpdfstring{\(\phi_{0,1}\)}{phi\_{0,1}}}
\label{sec:singular-theta-lift}

The automorphic, Borcherds--Kac--Moody, and singular-theta descriptions
of \(\Delta_5\) are theorems of Borcherds, Gritsenko, and Nikulin.
The Pfaffian description is relative to the retained Dirac--Igusa data,
and the mirror-discriminant description is conjectural.  The
exterior-square isogeny
\(\PSp_4(\mathbb R)\simeq\SO(3,2)_+\) identifies the genus-two Siegel
upper half-space and the type-IV symmetric domain attached to the
signature-\((3,2)\) lattice into one homogeneous space; the singular
theta integral of \cite{Borcherds95,B} sends the Weil-representation
theta component \(\phi^{\mathrm{Weil}}_{0,1}\) of \(\phi_{0,1}\) to
the weight-five Borcherds product \(\Delta_5\), and the
Gritsenko--Nikulin denominator identity then reads the lift as the BKM denominator on
\(\La^{2,1}_{II}\).  These identifications are theorems of
\cite{Borcherds95,B,GN,GNII}; the integral lattice, arithmetic
subgroup, and multiplier convention are fixed by the following
normalization.

\subsection{Notation and the canonical name}
\label{ssec:theta-lift-notation}

Three Borcherds-product handles for the weight-five Igusa cusp form
are distinguished.  The \emph{automorphic section}
\(\mathcal D_X=\Delta_5\) of the line
\(\mathcal L^5\otimes\nu_{\Delta_5}\) on
\(\Sp_4(\Z)\backslash\Hpl_2\) is the View-I name.  The \emph{BKM
denominator}
\(\operatorname{den}(\autcor)=64^{-1}\Delta_5(2Z)\)
on \(\La^{2,1}_{II}\) is the View-II name.  The \emph{protected
Pfaffian} \(\Pf_{\mathrm{prot}}(\mathfrak D_X^{\mathrm{DI}})\), defined on
\(K3\times E\) relative to the retained D0-HN datum, retained quotient-orientation
datum, chosen Weyl lifts, \((O2_{\mathrm{atlas}})\), and the finite
geometric Pfaffian datum \((P_{\mathrm{fin}})\), is the
View-III name.  The View-V canonical name is
\[
\Delta_5
\;=\;
\Theta\bigl(\phi^{\mathrm{Weil}}_{0,1}\bigr),
\]
the \emph{singular theta lift of the Weil component
\(\phi^{\mathrm{Weil}}_{0,1}\)} obtained from \(\phi_{0,1}\) by theta
decomposition.  Throughout
\(\phi_{0,1}\in J_{0,1}^{\mathrm{weak}}\) is the unique up to scale
weight-zero index-one weak Jacobi form on the Jacobi group
\(\Gamma^J_1=\SL_2(\Z)\ltimes\Z^2\) with constant Fourier coefficient
fixed by the K3 elliptic genus normalization
\(Z_{\mathrm{K3}}=2\phi_{0,1}\) of \cite[\S2.2]{EZ:1};
\(\phi^{\mathrm{Weil}}_{0,1}\) is its weight-\(-1/2\)
Weil-representation theta component, and \(\Theta\) is the Borcherds
singular-theta integral of \cite[\S6]{Borcherds95}, \cite[\S6]{B}
with the Igusa cusp trivialization
\(C_{\phi_{0,1}}=64\).  Its monic Borcherds product is
\(D_5=64^{-1}\Delta_5\).  The bare handle \(\Delta_5\) is reserved for
cross-view identifications.

\subsection{The weak Jacobi form \texorpdfstring{\(\phi_{0,1}\)}{phi\_{0,1}}}
\label{ssec:phi01}

\begin{definition}[Weak Jacobi form, weight zero, index one]
\label{def:weak-jacobi-phi01}
A weak Jacobi form of weight \(k\) and index \(m\) on the Jacobi group
\(\Gamma^J_1=\SL_2(\Z)\ltimes\Z^2\) is a holomorphic function
\(\phi:\Hpl\times\C\rightarrow\C\) satisfying, for
\(\begin{psmallmatrix}a&b\\c&d\end{psmallmatrix}\in\SL_2(\Z)\) and \((\lambda,\mu)\in\Z^2\),
\begin{align*}
\phi\!\left(\frac{a\tau+b}{c\tau+d},\ \frac{z}{c\tau+d}\right)
&=(c\tau+d)^k e^{2\pi i mcz^2/(c\tau+d)}\phi(\tau,z),\\
\phi(\tau,\,z+\lambda\tau+\mu)
&=e^{-2\pi i m(\lambda^2\tau+2\lambda z)}\phi(\tau,z),
\end{align*}
together with the cusp expansion
\[
\phi(\tau,z)=\sum_{n\geq0,\,l\in\Z}f(n,l)\,q^nr^l,
\qquad q=e^{2\pi i\tau},\ r=e^{2\pi iz},
\]
in which the coefficient depends only on the pair
\((4mn-l^2,\;l\bmod 2m)\) (Eichler--Zagier
\cite[Theorem 2.2]{EZ:1}).  For \(m=1\) the residue class is fixed by
the parity of \(4n-l^2\).  The
\emph{weight-zero index-one} weak Jacobi form \(\phi_{0,1}\) is fixed
by \(k=0\), \(m=1\), the discriminant condition
\(4n-l^2\geq -1\) and the normalization
\[
\phi_{0,1}(\tau,z)=r^{-1}+10+r+O(q).
\]
\end{definition}

The space \(J_{0,1}^{\mathrm{weak}}\) is one-dimensional over \(\C\)
(Eichler--Zagier \cite[Theorem 9.4 and Table 9, \S9]{EZ:1});
\(\phi_{0,1}\) generates it.  The theta-quotient presentation is
the input to the singular-theta integral below.

\begin{lemma}[Eichler--Zagier theta-quotient presentation]
\label{lem:phi01-presentations}
The weak Jacobi form \(\phi_{0,1}\) admits the closed form
\[
\phi_{0,1}(\tau,z)
=4\sum_{i=2}^{4}
\frac{\theta_i(\tau,z)^2}{\theta_i(\tau,0)^2},
\]
where \(\theta_2,\theta_3,\theta_4\) are the three even Jacobi theta
functions in the Mumford normalization (Eichler--Zagier
\cite[\S9, equation (9.4)]{EZ:1}); equivalently
\(\phi_{0,1}=Z_{\mathrm{K3}}/2\) is one half of the K3 elliptic genus
\cite{DMVV}.  The leading Fourier coefficients are
\[
f(0,0)=10,\quad f(0,\pm1)=1,\quad f(1,0)=108,\quad
f(1,\pm1)=-64,\quad f(1,\pm2)=10,
\]
\[
f(2,0)=808,\quad f(2,\pm1)=-513,\quad f(2,\pm2)=108,\quad f(2,\pm3)=1,
\]
in agreement with the Mukai--Gram normalization of
Appendix~\ref{appendix:mukai-gram-dictionary}.
The discriminant relation \(f(n,l)=f(n',l')\) when
\(4n-l^2=4n'-(l')^2\) (Eichler--Zagier \cite[Theorem~2.2]{EZ:1})
is the index-one specialization of the
\((4mn-l^2,l\bmod 2m)\) rule and organises these coefficients into
the equivalence classes
\[
\{f(1,0),\,f(2,\pm2)\}=108\ (\text{disc }4),\quad
\{f(0,0),\,f(1,\pm2)\}=10\ (\text{disc }0),\quad
\{f(0,\pm1),\,f(2,\pm3)\}=1\ (\text{disc }-1),
\]
each independently consistent with the verified expansion.
The finite arithmetic table
\[
\texttt{certificates/jacobi/k3e\_first\_relation\_window\_coefficients}
\]
records the computed rows \(\phi_{0,1}=r^{-1}+10+r+O(q)\), the finite
evenness checks \(f(n,l)=f(n,-l)\), and the finite
discriminant-dependence checks \(f(n,l)=c(4n-l^2)\) for the first
relation-closed window.  This table checks the displayed Fourier
window; the all-coefficient statement is the Eichler--Zagier theorem
cited above.
\end{lemma}

\begin{proof}
The closed form is standard in Eichler--Zagier
\cite[\S\S9--10]{EZ:1}; the K3 elliptic-genus identity
\(Z_{\mathrm{K3}}=2\phi_{0,1}\) is \cite[\S2.2]{EZ:1} together with the
sigma-model right-moving Witten index of \cite{DMVV}.  The constant
term \(f(0,0)=10\) is the right-moving-localized character of the K3
Ramond sector at \(L_0=c/24\); the remaining values follow from the
discriminant relation \(f(n,l)=f(n',l')\) when
\(4n-l^2=4n'-(l')^2\) (Eichler--Zagier \cite[Theorem 2.2]{EZ:1})
together with explicit theta-quotient evaluation, which fixes the
disc-\(3\) class (\(f(1,\pm1)=-64\)) and the disc-\(4\) class
(\(f(1,0)=f(2,\pm2)=108\)) by direct computation from the
theta-quotient of Lemma~\ref{lem:phi01-presentations}.  The squared
theta-leading coefficient
\(([q^{1/2}r^{1/2}s^{1/2}]\Delta_5)^2 = 64^2 = 4096\) of the
Borcherds-product expansion in
\S\ref{appendix:mukai-gram-dictionary} is the next
arithmetic check on these coefficients via the cusp-chamber product
\(\Delta_5(Z) = 64q^{1/2}r^{1/2}s^{1/2}\prod_{(n,l,m)\in\Gamma_{\mathrm{eff}}}
(1-q^nr^ls^m)^{f(nm,l)}\).
\end{proof}

\subsection{The signature-\texorpdfstring{\((3,2)\)}{(3,2)} lattice and the type-IV domain}
\label{ssec:lattice-3-2}

\begin{definition}[The signature-\((3,2)\) lattice]
\label{def:lambda-3-2}
The lattice \(\La^{3,2}\) is the orthogonal complement of the alternating
Pfaffian element
\(q_0=e_1\wedge e_3+e_2\wedge e_4\in\wedge^2\La^4\) inside the
signature-\((3,3)\) lattice
\(\bigl(\wedge^2\La^4,\,(\,,\,)\bigr)\simeq\mathrm{II}_{3,3}\),
where \(\La^4=\bigoplus_{i=1}^4\Z e_i\) and the bilinear form
\((u,v)\) is the Pfaffian pairing
\(u\wedge v=(u,v)\,e_1\wedge e_2\wedge e_3\wedge e_4\).  The vector
\(q_0\) has square \(-2\).  In the BKM sign convention it is the even discriminant-two lattice
\[
\La^{3,2}\;\simeq\;\HypPlan\oplus\HypPlan\oplus[2],
\qquad
\HypPlan=\begin{psmallmatrix}0&-1\\-1&0\end{psmallmatrix},
\]
written in the basis
\((f_1,f_2,f_3,f_{-2},f_{-1})\) of
the exterior-square model section.
\end{definition}

\begin{remark}[Signature conventions]
The lattice \(\La^{3,2}\) carries three positive and two negative
directions in the \((u,v)\)-pairing of
Definition~\ref{def:lambda-3-2}.  The \emph{primitive hyperbolic
sublattice} \(\La^{2,1}=\HypPlan\oplus[2]\) on which the BKM denominator
algebra \(\autcor\) is graded has signature \((2,1)\) in the fixed
convention: the \([2]\) factor and one hyperbolic
direction are positive, and the other hyperbolic direction is negative.
The decomposition
\(\La^{3,2}=\La^{2,1}\oplus\HypPlan\) splits an extra hyperbolic plane
on which the singular theta integral lives; this is the
\(\mathrm{II}_{2,2}\)-fibre at the standard cusp under the exceptional
isomorphism of \S\ref{ssec:exceptional-iso}.  No
\(\La^{2,1}\oplus\langle-2\rangle\) presentation is invoked; the
fixed convention is the \(\HypPlan\oplus[2]\) Nikulin form.
\end{remark}

Let
\[
\Gamma_{3,2}:=\wedge^2\Sp_4(\Z)/\{\pm\Id_4\}
\subset \mathrm P\Orth(\La^{3,2})_+
\]
be the projective exterior-square arithmetic group.  The type-IV
symmetric domain attached to \(\La^{3,2}\) is
\[
\Hpl^{\IV}=\bigl\{Z\in\Proj(\La^{3,2}\otimes\C)\,\bigm|\,
(Z,Z)=0,\ (Z,\overline Z)<0\bigr\},
\]
with positive component \(\Hpl^{\IV}_+\) singled out by
\(\Imag z_1>0\) in the parametrization
\({}^t(z_2^2-z_1z_3,z_3,z_2,z_1,1)z_0\) of
the exterior-square model section.  Its arithmetic quotient
\(\Gamma_{3,2}\backslash\Hpl^{\IV}_+\) is the moduli space on
which the Borcherds singular-theta lift of \cite{Borcherds95} produces
automorphic forms with prescribed singularities along rational quadratic
divisors.

\subsection{The exceptional isogeny \texorpdfstring{\(\PSp_4\simeq\SO(3,2)_+\)}{PSp(4) ~= SO(3,2)+}}
\label{ssec:exceptional-iso}

\begin{theorem}[Exterior-square exceptional isomorphism]
\label{thm:exceptional-iso}
The exterior-square representation
\(\wedge^2:\SL_4\rightarrow\Aut(\wedge^2\La^4,(\,,\,))\) restricts to a
homomorphism
\[
\wedge^2:\Sp_4(\Z)/\{\pm\Id_4\}\;\xrightarrow{\ \simeq\ }\;\Gamma_{3,2},
\]
realizing the exceptional isomorphism of adjoint real Lie groups
\(\PSp_4(\R)\simeq\SO(3,2)_+\), together with its spin lift on the
covering groups.  Under this identification the genus-two
Siegel upper half-space \(\Hpl_2\) embeds onto the positive component
\(\Hpl^{\IV}_+\) by
\[
\begin{psmallmatrix}z_1&z_2\\z_2&z_3\end{psmallmatrix}\;\longmapsto\;
{}^t(z_2^2-z_1z_3,\ z_3,\ z_2,\ z_1,\ 1)z_0,
\]
intertwining the \(\Sp_4(\Z)\) and \(\Gamma_{3,2}\) actions.
\end{theorem}

\begin{proof}
The Pfaffian form
\(u\wedge v=(u,v)e_1\wedge e_2\wedge e_3\wedge e_4\) is the canonical
\(\SL_4\)-invariant symmetric bilinear form on \(\wedge^2\La^4\); the
subgroup of \(\SL_4\) fixing \(q_0=e_1\wedge e_3+e_2\wedge e_4\) is the
symplectic group \(\Sp_4\) of the alternating form \(B_{q_0}\).
The orthogonal complement \(\La^{3,2}=q_0^\perp\) is preserved, and the
kernel of \(\wedge^2|_{\Sp_4}\) on \(q_0^\perp\) is \(\{\pm\Id_4\}\)
(the exterior-square model section).  This is the integral form of
the classical exceptional isomorphism \(\PSp_4(\R)\simeq\SO(3,2)_+\)
of \cite[Chapter II.6]{KNAPP:1}; the spin cover supplies the
corresponding lift with Lie algebra \(\mathfrak{sp}_4(\R)\).  This cover is distinct
from the metaplectic cover on the \(\Mp_2(\Z)\) input.  The displayed
substitution into the type-IV domain produces the commutative square of
the exterior-square model section.  The integral-Igusa form of the
identification \(\Hpl_2\simeq\Hpl^{\IV}_+\) is the content of
\cite[\S2]{Igusa1962}; see also \cite[Chapter II]{FREITAG:1} and
\cite[\S6]{VDGEER:1}.
\end{proof}

\begin{remark}[Three guises of one homogeneous space]
The same homogeneous space carries three names: the genus-two Siegel
domain \(\Sp_4(\R)/\Un(2)\), the type-IV symmetric domain
\(\SO(3,2)_+/(\SO(3)\times\SO(2))\), and the spin double cover
\(\Spin(3,2)/(\Un(2)\text{-cover})\).  The regularized theta integral
below uses the metaplectic cover \(\Mp_2(\Z)\) on the input variable
\(\tau\) through the Weil representation of the discriminant form of
\(\La^{3,2}\).  The integrality of \(\operatorname{wt}(\Delta_5)=5\)
places the resulting weight on the ordinary automorphic line; the multiplier system
\(\nu_{\Delta_5}\) records the residual reflection character on
\(W^{(2)}(\La^{2,1}_{II})\) (Lemma~\ref{lem:bkm-maass-restriction}).
\end{remark}

\subsection{The Borcherds theta correspondence and the exceptional isomorphism}
\label{ssec:borcherds-theta-correspondence}

\begin{theorem}[Regularized theta correspondence on \(\La^{3,2}\)]
\label{thm:regularized-theta-correspondence}
Let \(L=\La^{3,2}\).  Borcherds's regularized theta integral pairs a
weakly holomorphic form
\(\phi\in M^!_{-1/2}(\rho_L)\) for \(\Mp_2(\Z)\) with the
non-holomorphic Siegel theta kernel of \(L\), and produces an
automorphic form on \(\widetilde\Orth(L)_+\) of weight
\(c_\phi(0,0)/2\).  Its rational-quadratic divisor is determined by the
negative-index part of \(\phi\), and its Weyl-chamber product uses the
Fourier coefficients of \(\phi\) in all effective product degrees.  For
\(\phi=\phi^{\mathrm{Weil}}_{0,1}\), the pullback along
\[
\wedge^2:\Sp_4(\Z)/\{\pm\Id_4\}\longrightarrow\Gamma_{3,2}
\subset\Orth(\La^{3,2})
\]
is the Igusa cusp form \(\Delta_5\).
\end{theorem}
\hfill\textnormal{[proved; \cite{Borcherds95,B,Igusa1962,GN}]}

\begin{remark}[The two groups in the lift]
The metaplectic group in the construction is \(\Mp_2(\Z)\), acting on
the input form through the Weil representation \(\rho_{\La^{3,2}}\).
The genus-two group \(\Sp_4\) enters after the orthogonal automorphic
form has been constructed, through the exceptional isomorphism
\(\PSp_4(\R)\simeq\SO(3,2)_+\) and the integral exterior-square map of
Theorem~\ref{thm:exceptional-iso}.  Thus the theta input, the
orthogonal automorphic form, and the Siegel modular form are three
successive incarnations of one Borcherds product.
\end{remark}

\subsection{The Borcherds singular-theta integral}
\label{ssec:singular-theta}

Borcherds's correspondence assigns to a Weil-representation form
obtained from a Jacobi input a meromorphic orthogonal modular form by
integration against a vector-valued Siegel theta series.  The integral
is divergent at the lift's input weight; the Borcherds regularization
produces a meromorphic Siegel modular form, after the
\(\PSp_4\simeq\SO(3,2)_+\) pullback, with explicit divisor.

\begin{definition}[Siegel theta-series of \texorpdfstring{\(\La^{3,2}\)}{Lambda^{3,2}}]
\label{def:siegel-theta}
For the lattice \(\La^{3,2}\) of Definition~\ref{def:lambda-3-2}, the
Siegel theta-series at \(Z\in\Hpl^{\IV}_+\) and \(\tau\in\Hpl\) is
\[
\Theta_{\La^{3,2}}(\tau,Z)
=\sum_{\gamma\in A_{\La^{3,2}}}e_\gamma
\sum_{\lambda\in\La^{3,2}+\gamma}
\exp\!\bigl(\pi i\,\overline\tau\,(\lambda,\lambda)_{Z^+}
+\pi i\,\tau\,(\lambda,\lambda)_{Z^-}\bigr),
\]
where \((\,,\,)_{Z^\pm}\) are the positive and negative parts of the
\((u,v)\)-pairing under the orthogonal decomposition determined by
\(Z\) (Borcherds \cite[\S4]{Borcherds95}, \cite[\S4]{B}).  The
theta-series is a non-holomorphic modular form of weight
\((3/2,1)\) on \(\Mp_2(\Z)\) valued in the Weil representation of
\(\La^{3,2}\) on \(\C[A_{\La^{3,2}}]\).  The scalar sum over
\(\La^{3,2}\) is only the \(\gamma=0\) component.
\end{definition}

\begin{theorem}[Singular theta lift {\cite[Theorem 13.3]{Borcherds95}, \cite[Theorem 6.2]{B}}]
\label{thm:singular-theta-lift}
Let \(\phi\) be a weakly holomorphic vector-valued modular form for the
dual Weil representation of \(\La^{3,2}\) on \(\Mp_2(\Z)\) of weight
\(-1/2\) with integral Fourier coefficients
\(c_\phi(n,\lambda)\in\Z\) at the cusp.  The regularized integral
\[
\Theta(\phi)(Z)
=\int^{\mathrm{reg}}_{\SL_2(\Z)\backslash\Hpl}
\bigl\langle\phi(\tau),\Theta_{\La^{3,2}}(\tau,Z)\bigr\rangle\,v
\frac{du\,dv}{v^2},
\qquad\tau=u+iv,
\]
in the regularization sense of \cite[\S6]{B} produces a meromorphic
function on \(\Hpl^{\IV}_+\) with the following properties.
\begin{enumerate}[label=\textup{(\arabic*)}]
\item \emph{Modular invariance.} \(\Theta(\phi)\) is an automorphic
form of weight \(c_\phi(0,0)/2\) on
\(\widetilde\Orth(\La^{3,2})_+\) with multiplier system determined by
the Weil action on \(\La^{3,2}{}^*/\La^{3,2}\).
\item \emph{Borcherds product.} On the cone of definition,
\(\Theta(\phi)\) admits the Borcherds product expansion
\[
\Theta(\phi)(Z)
=C_\phi e^{-2\pi i(\rho,Z)}
\prod_{\substack{\alpha\in\La^{3,2}{}^*\\(\alpha,W)>0}}
	\bigl(1-e^{-2\pi i(\alpha,Z)}\bigr)^{c_\phi(-(\alpha,\alpha)/2,\,\alpha\bmod\La^{3,2})},
\]
with \(C_\phi\in\C^\times\) fixed by the chosen cusp trivialization,
Weyl vector \(\rho\) determined by \(\phi\), and a chamber \(W\) of the
positive cone.
\item \emph{Divisor.} The divisor of \(\Theta(\phi)\) is the linear
combination of rational quadratic divisors \(\lambda^\perp\)
weighted by \(c_\phi(-(\lambda,\lambda)/2,\lambda\bmod\La^{3,2})\)
over primitive \(\lambda\in\La^{3,2}{}^*\) with negative discriminant.
\end{enumerate}
The factor \(v\) is the invariant normalization for signature
\((3,2)\): the input has holomorphic weight \(-1/2\), the theta kernel
has weight \((3/2,1)\), and
\(v\langle\phi,\Theta_{\La^{3,2}}\rangle\) has weight \((0,0)\).
\end{theorem}
\hfill\textnormal{[proved; \cite{Borcherds95,B}]}

\begin{remark}[Two regularization conventions]
The Borcherds regularization \cite[\S6]{B} subtracts the divergent
\(v\rightarrow\infty\) contribution by analytic continuation in an
auxiliary parameter, and is the convention used here.  The
Harvey--Moore regularization \cite{HM} differs from the Borcherds
regularization by a finite renormalization of the Weyl vector and the
constant term, with no effect on the product expansion's exponents or
divisor.  The two conventions agree on quotient divisors
\(\Theta(\phi_1)/\Theta(\phi_2)\), on the BKM denominator identity, and on the weight identity
\(\operatorname{wt}\Theta(\phi)=c_\phi(0,0)/2\).
\end{remark}

\subsection{From \texorpdfstring{\(\phi_{0,1}\)}{phi\_{0,1}} to a vector-valued Weil-representation form}
\label{ssec:phi-to-weil}

The singular-theta integral takes a weakly holomorphic vector-valued
form on \(\Mp_2(\Z)\); the input \(\phi_{0,1}\) is a scalar weak Jacobi
form of index one.  Theta-decomposition reduces the latter to the
former.

\begin{lemma}[Theta-decomposition of \texorpdfstring{\(\phi_{0,1}\)}{phi\_{0,1}}]
\label{lem:theta-decomposition}
Let \(\phi_{0,1}\in J_{0,1}^{\mathrm{weak}}\) and decompose along the
elliptic theta-functions of index one,
\[
\phi_{0,1}(\tau,z)=h_0(\tau)\theta_{1,0}(\tau,z)+h_1(\tau)\theta_{1,1}(\tau,z),
\]
where \(\theta_{1,\mu}(\tau,z)=\sum_{\ell\equiv\mu\, (2)}
q^{\ell^2/4}r^\ell\) for
\(\mu\in\{0,1\}\) (Eichler--Zagier \cite[\S5]{EZ:1}).  Then
\(h=(h_0,h_1)\) is a weakly holomorphic vector-valued modular form of
weight \(-1/2\) for the dual Weil representation
\(\rho_{[2]}^\vee\), equivalently for the Weil representation of the
opposite discriminant form \([-2]\), on \(\Mp_2(\Z)\).  Under the
finite quadratic-module identification
\[
A_{[2]}\;\simeq\;A_{\HypPlan\oplus[2]}
\;\simeq\;A_{\La^{3,2}},
\]
it pairs with the \(\rho_{\La^{3,2}}\)-valued theta kernel of
Definition~\ref{def:siegel-theta}.  Its Fourier coefficients are exactly
\[
[q^{D/4}]h_\mu(\tau)=c(D),\qquad
c(4n-l^2)=f(n,l),\qquad D\equiv-\mu^2\pmod4,
\]
or equivalently \([q^{\,n-l^2/4}]h_{l\bmod 2}=f(n,l)\).  In the
index-one normalization the first coefficients are
\([q^0]h_0=10\) and \([q^{-1/4}]h_1=1\) reproducing
\(\phi_{0,1}=r^{-1}+10+r+O(q)\).  Since
\(\theta_{1,1}(\tau+1,z)=i\,\theta_{1,1}(\tau,z)\), invariance of
\(\phi_{0,1}\) under \(T\) forces
\(h_1(\tau+1)=-i\,h_1(\tau)\), the dual eigenvalue.  The vector \(h\)
is the weakly holomorphic input
\(\phi^{\mathrm{Weil}}_{0,1}\) of the weight and representation type
required by Theorem~\ref{thm:singular-theta-lift}.
\end{lemma}
\hfill\textnormal{[proved; \cite{EZ:1,Borcherds95,GN}]}

\begin{proof}
Theta-decomposition along the lattice \(\La_{\mathrm{ind}}=\Z\) inside
the Jacobi-elliptic variable is Eichler--Zagier
\cite[Theorem 5.1]{EZ:1}.  The two-component vector \(h\) is weakly
holomorphic of weight \(-1/2\) on \(\Mp_2(\Z)\) with the Weil
representation of the rank-one lattice \(\langle 2\rangle\).  Comparing
the coefficient of \(r^\ell\) gives
\([q^{\,n-\ell^2/4}]h_{\ell\bmod2}=f(n,\ell)\), hence the discriminant
formula above.  The hyperbolic plane \(\HypPlan\) is unimodular, so
adjoining one or two copies of \(\HypPlan\) does not change the
discriminant form.  Thus
\[
(\HypPlan\oplus[2])^*/(\HypPlan\oplus[2])
\simeq
[2]^*/[2]
\simeq
(\La^{3,2})^*/\La^{3,2},
\]
with the same finite quadratic form.  This discriminant-form isometry,
is the transport of \(h\) to the required weakly holomorphic
vector-valued form for the Weil representation of \(\La^{3,2}\).
The constant Fourier coefficient
\(c_{\phi^{\mathrm{Weil}}_{0,1}}(0,0)=f(0,0)=10\) is preserved under
this discriminant-form identification because the trivial coset
\(0\in\La^{3,2}{}^*/\La^{3,2}\) matches the trivial coset on the
rank-one factor.  Borcherds
\cite[\S15]{Borcherds95} records the explicit computation for the
\(\Sp_4\)-case in the equivalent paramodular language.
\end{proof}

\begin{remark}[Index-one normalization of \(c_1(0)\)]
The constant Fourier coefficient
\(c_1(0):=f(0,0)=10\) of \(\phi_{0,1}\) is the Eichler--Zagier
index-one normalization of the K3 elliptic genus.  The
universal identity \(\kappa_{\mathrm{BKM}}(\Phi_N)=c_N(0)/2\) of
\cite{Gritsenko99} names \(N\) as the paramodular level of the input
denominator;
at \(N=1\) the input form is \(\phi_{0,1}\) and \(c_1(0)=10\), giving
\[
\kappa_{\mathrm{BKM}}(\Delta_5)=c_1(0)/2=5.
\]
The integer-versus-half-integer ambiguity for the metaplectic weight
factor is killed by the integrality of \(c_1(0)/2=5\), and the
Eichler--Zagier normalization is the row \(F_1\) normalization in the
Cl\'ery--Gritsenko catalogue \cite{GC}.  On the other seven rows the
symbol \(c(0)\) is indexed by the pair \((t,N)\), not by a CHL twining
order alone.  No \(c_N(0)=2k_N+1\) half-integer convention enters at
\(N=1\).
\end{remark}

\subsection{Theta lift identifies \texorpdfstring{\(\Delta_5\)}{Delta\_5}}
\label{ssec:theta-lift-delta5}

The lift produces an automorphic form on
\(\widetilde\Orth(\La^{3,2})_+\), pulled back along
\(\wedge^2:\Sp_4(\Z)/\{\pm\Id_4\}\xrightarrow{\simeq}\Gamma_{3,2}\)
of Theorem~\ref{thm:exceptional-iso} to a Siegel modular form on
\(\Sp_4(\Z)\).

\begin{theorem}[\texorpdfstring{\(\Delta_5\)}{Delta\_5} as singular theta lift]
\label{thm:delta5-as-theta-lift}
The singular-theta lift of the weak Jacobi form \(\phi_{0,1}\) is the
weight-five Igusa cusp form:
\[
\Theta\bigl(\phi^{\mathrm{Weil}}_{0,1}\bigr)\!\bigm|_{\Hpl_2}
\;=\;\Delta_5,
\]
the right-hand side written in the convention
\(\mathcal D_X=\Delta_5\) of \S\ref{ssec:theta-lift-notation} and the
left-hand side regularized in the Borcherds sense of
\cite[\S6]{B}.  Equivalently, in the BKM normalization,
\[
\Theta\bigl(\phi^{\mathrm{Weil}}_{0,1}\bigr)
=64\cdot\operatorname{den}(\autcor)\!\bigm|_{Z\mapsto Z/2},
\]
i.e.\ the lift is \(\Delta_5\) on \(\Hpl_2\) and
\(64^{-1}\Delta_5(2Z)\) on \(\La^{2,1}_{II}\)-coordinates of
\(\Hpl^{\IV}_+\).
\end{theorem}
\hfill\textnormal{[proved; \cite{Borcherds95,B,GN,GNII}]}

\begin{proof}
The Borcherds singular-theta lift in
Theorem~\ref{thm:singular-theta-lift} produces a meromorphic
automorphic form
\(\Theta(\phi^{\mathrm{Weil}}_{0,1})\) on
\(\widetilde\Orth(\La^{3,2})_+\) of weight \(c_1(0)/2=5\).  Its
divisor is the linear combination of rational quadratic divisors
weighted by \(c_{\phi^{\mathrm{Weil}}_{0,1}}(-(\lambda,\lambda)/2,
\lambda\bmod\La^{3,2})\).  For \(\phi_{0,1}\), the
negative-discriminant coefficients \(f(0,\pm1)=1\) give the simple
reflective Humbert divisor of \(\Delta_5\), while \(f(0,0)=10\) gives
the weight \(5\).  The finite packet
\(\texttt{certificates/automorphic/delta5\_humbert\_divisor}\)
records this divisor support and order-one multiplicity separately from
the pole-order-two statement for \(\Delta_5^{-2}\).  The
Borcherds-product expansion of
Theorem~\ref{thm:singular-theta-lift}(2), specialized to
\(\phi^{\mathrm{Weil}}_{0,1}\), reads
\[
\Theta(\phi^{\mathrm{Weil}}_{0,1})
=C_{\phi_{0,1}}q^{1/2}r^{1/2}s^{1/2}
\prod_{(n,l,m)\in\Gamma_{\mathrm{eff}}}
(1-q^nr^ls^m)^{f(nm,l)}
\]
in the cusp chamber, with \(\Gamma_{\mathrm{eff}}\) the cusp-positive
semigroup and Weyl vector
\(\rho_0=(1/2,1/2,1/2)\).  The monic Borcherds product is
\(D_5(Z)\) of \S\ref{ssec:theta-lift-notation}; the Igusa
theta-constant cusp trivialization fixes
\(C_{\phi_{0,1}}=64\), hence
\(\Theta(\phi^{\mathrm{Weil}}_{0,1})=\Delta_5\).  The pullback to \(\Hpl_2\) along
\(\wedge^2:\Sp_4(\Z)/\{\pm\Id_4\}\rightarrow\Gamma_{3,2}\) of
Theorem~\ref{thm:exceptional-iso} identifies
\(\Theta(\phi^{\mathrm{Weil}}_{0,1})\!\bigm|_{\Hpl_2}\) with the
\(\Sp_4(\Z)\)-cusp form whose Borcherds product is exactly the Igusa
cusp form \(\Delta_5\) of weight five.  The BKM normalization
\(64^{-1}\Delta_5(2Z)\) on \(\La^{2,1}_{II}\) is
Gritsenko--Nikulin \cite[Theorem 2.1]{GN}, \cite[Theorem 2.1, Section
5.1.1]{GNII}: the doubled lattice argument \(Z\mapsto 2Z\) corresponds
to the polarization change \(\alpha:\Gamma_{\mathrm{ind}}
\rightarrow\La^{2,1}_{II}\) of \S\ref{ssec:weyl-vector-rho}.  Both
sides have the same Borcherds-product expansion, and their leading
coefficients agree by the theta-constant
\([q^{1/2}r^{1/2}s^{1/2}]\Delta_5=64\).
\end{proof}

\begin{corollary}[Borcherds--Gritsenko--Nikulin pentadic identity]
\label{cor:five-views-coincidence}
The theorem-level automorphic, BKM-denominator, and theta-lift readings
coincide:
\begin{align*}
\mathrm{(View\ I)}\ \mathcal D_X
&=\mathrm{(View\ II)}\ 64\cdot\operatorname{den}(\autcor)\!\bigm|_{Z\mapsto Z/2}\\
&=\mathrm{(View\ V)}\ \Theta\bigl(\phi^{\mathrm{Weil}}_{0,1}\bigr).
\end{align*}
The Pfaffian identification
\(\Pf_{\mathrm{prot}}(\mathfrak D_X^{\mathrm{DI}})=\Delta_5\) of View III
holds on \(K3\times E\) at the Pfaffian and orientation level, granted the
retained D0-HN datum, retained quotient-orientation datum, chosen Weyl lifts,
retained \((O2_{\mathrm{atlas}})\) wall datum, and retained finite
geometric Pfaffian datum \((P_{\mathrm{fin}})\) of
Theorem~\ref{thm:G-four-obstruction-criterion} of
Appendix~\ref{app:obstruction-discharges}.  The lift does not
provide the level-\(\mathsf A\)
gravity-line operator algebra; a 3d-gravity path-integral realization
requires the level-\(\mathsf A\) Hall--Borcherds residual, while
Definition~\ref{def:G-vol2-hall-borcherds} records only the
level-\(\mathsf Z\) trace comparison.
\end{corollary}
\hfill\textnormal{[proved for Views I, II, V; conditional for View III]}

\subsection{The multiplier system and the Weyl vector}
\label{ssec:weyl-vector-rho}

\begin{theorem}[Multiplier system of the lift]
\label{thm:multiplier-system}
The lift \(\Theta(\phi^{\mathrm{Weil}}_{0,1})\) carries the
\(\Sp_4(\Z)\)-multiplier system \(\nu_{\Delta_5}\) of Maass
\cite{MAASS:1}, equivalently described as the character
\(v_{\eta}^{12}\times v_H\) on the Jacobi parabolic and as the
Gritsenko character on the BKM side
\cite[Proposition 2.1]{GN}, \cite[\S5.1.1, case \((t=1,\mathrm{II},0)\)]{GNII}.
On the type-II generators of the Weyl group
\(W^{(2)}(\La^{2,1}_{II})\),
\[
\nu_{\Delta_5}(s_{\delta_i})=-1,\qquad i=1,2,3,
\]
and on the chamber automorphism factor \(\Aut(\Poly_{II})\simeq S_3\)
the multiplier is trivial:
\[
\nu_{\Delta_5}(s_{f_{-2}-f_2})
=\nu_{\Delta_5}(s_{f_2-f_3})
=\nu_{\Delta_5}(s_{f_{-2}-f_3})=+1.
\]
This is exactly the multiplier of
Lemma~\ref{lem:bkm-maass-restriction}.
The finite automorphic packet
\(\texttt{certificates/automorphic/delta5\_maass\_character}\)
checks these generator values, the order-two relation, and the
automorphic-line placement \( \mathcal L^5\otimes\nu_{\Delta_5} \);
the packet is not an orientation construction.
\end{theorem}
\hfill\textnormal{[proved; \cite{MAASS:1,GN,GNII}]}

\begin{remark}[Singular-theta multiplier vs Maass multiplier]
The Borcherds singular-theta lift produces a multiplier system on
\(\widetilde\Orth(\La^{3,2})_+\) determined by the Weil action on
\(\La^{3,2}{}^*/\La^{3,2}\) (Borcherds \cite[Theorem 13.3]{Borcherds95}).
For the present \(\La^{3,2}\), the discriminant group has order two,
generated by \(f_3/2\) in the convention of
Remark~\ref{rem:discriminant-conventions}.  The finite Weil module is
the two-dimensional module attached to this discriminant form.  The
Borcherds multiplier determined by that module and by the Weyl vector
restricts on the Borcherds chamber to the same reflection character.
Pulled back along
\(\wedge^2:\Sp_4(\Z)/\{\pm\Id_4\}\rightarrow\Gamma_{3,2}\) it
matches Maass's multiplier formula for the product of the ten even
theta constants \cite[(1.3)--(1.5)]{GN}, \cite{MAASS:1}: the
\(-1\)-eigenvalues on the three type-II simple reflections and the
\(+1\)-eigenvalues on the three chamber automorphisms.  The two
descriptions are not independent — they are the two sides of the same
character — and their agreement is part of the regularized theta
correspondence of Theorem~\ref{thm:regularized-theta-correspondence}.
\end{remark}

\begin{theorem}[Weyl vector of the lift]
\label{thm:weyl-vector-lift}
The Weyl vector of the Borcherds product
\(\Theta(\phi^{\mathrm{Weil}}_{0,1})\) in the type-II chamber
\(\Poly_{II}\subset\La^{2,1}_{II}\otimes\R\) is
\[
\rho=\tfrac12\delta_1+\tfrac12\delta_2+\tfrac12\delta_3
=f_2-\tfrac12 f_3+f_{-2},
\]
satisfying \((\rho,\delta_i)=-1\) for \(i=1,2,3\), with simple-real-root
Cartan matrix
\[
A=\bigl((\delta_i,\delta_j)\bigr)
=\begin{pmatrix}2&-2&-2\\-2&2&-2\\-2&-2&2\end{pmatrix}.
\]
This is the Weyl vector of \(\autcor\) on \(\La^{2,1}_{II}\) of
the Weyl-chamber section, identified with the Weyl vector of
the Borcherds product by \cite[Theorem 4.1]{GN}.
The finite lattice packet
\(\texttt{certificates/lattice/type\_ii\_weyl\_vector}\) checks the
Cartan matrix, discriminant \(-32\), dual-lattice membership of
\(\rho\), and the displayed pairings.
\end{theorem}
\hfill\textnormal{[proved; \cite{GN,GNII}]}

\begin{proof}
The Weyl vector of this Borcherds product is determined by the
\(q^0\)-row of \(\phi\) at the cusp.  For
\(\phi_{0,1}=r^{-1}+10+r+O(q)\), the negative-discriminant coefficients
are \(f(0,\pm1)=1\), and the constant coefficient is \(f(0,0)=10\); the
Weyl-vector formula of \cite[Theorem 13.3, equation (13.10)]{Borcherds95}
yields \(\rho=(1/2,1/2,1/2)\) in cusp coordinates, transported to
\(\La^{2,1}_{II}\) by the polarization change
\(\alpha:\Gamma_{\mathrm{ind}}\rightarrow\La^{2,1}_{II}\) of
the Weyl-chamber section.  This recovers
\(\rho=f_2-\tfrac12 f_3+f_{-2}\); the Cartan matrix is the Gram matrix
of \(\{\delta_1,\delta_2,\delta_3\}\) computed in
the Weyl-chamber section.
\end{proof}

\subsection{Compatibility with View II: lift's image is the BKM denominator}
\label{ssec:compat-view-ii}

\begin{theorem}[Lift--denominator compatibility {\cite[Proposition 3.1]{GN}}]
\label{thm:lift-denominator}
The image of the singular-theta lift,
\(\Theta(\phi^{\mathrm{Weil}}_{0,1})=\Delta_5\), satisfies the
Weyl--Kac--Borcherds denominator identity for the Lorentzian
BKM superalgebra \(\autcor\) on \(\La^{2,1}_{II}\):
\[
64^{-1}\Delta_5(2Z)
=e^{-2\pi i(\rho,Z)}
\prod_{\alpha\in\mathcal R_+}
\bigl(1-e^{-2\pi i(\alpha,Z)}\bigr)^{\smult(\alpha)},
\]
with positive-root system \(\mathcal R_+\), Weyl vector \(\rho\) of
Theorem~\ref{thm:weyl-vector-lift}, and signed root supermultiplicities
\[
\smult\bigl(\alpha(n,l,m)\bigr)=f(nm,l),
\qquad (n,l,m)\in\Gamma_{\mathrm{act}}.
\]
The simple imaginary-root data are the Gritsenko--Nikulin coefficients
\(m(a)\) on negative-norm rays and \(\tau(a)\) on isotropic rays of
the Weyl-chamber section.
\end{theorem}
\hfill\textnormal{[proved; \cite{GN,GNII}]}

\begin{proof}
The Borcherds-product expansion
\(\Delta_5=64\,q^{1/2}r^{1/2}s^{1/2}
\prod(1-q^nr^ls^m)^{f(nm,l)}\) of
Theorem~\ref{thm:delta5-as-theta-lift} is the BKM denominator product
for \(\autcor\) after the polarization change
\(\alpha:\Gamma_{\mathrm{ind}}\rightarrow\La^{2,1}_{II}\): the exponent
\(f(nm,l)\) is the signed multiplicity of the root
\(\alpha(n,l,m)\), and the doubled argument \(2Z\) converts
\(\exp(-\pi i(\alpha,z))\) into the standard denominator character
\(\exp(-2\pi i(\alpha,z))\).
The finite packet
\(\texttt{certificates/lattice/product\_chamber\_root\_map}\)
checks these two target-side steps: the row-wise root map into
\(\La^{2,1}_{II}\) and the \(Z\mapsto2Z\) character conversion.
The algebra \(\autcor\) is the Gritsenko--Nikulin algebra of
\cite[\S\S3--4, Proposition 3.1]{GN}; the denominator identity is read
off the product expansion as in Chapter~\ref{ch:view2-bkm}.
Equivalently, the imaginary-simple-root characters \(m(a)\) and
\(\tau(a)\) are the additive data of
\cite[Theorem 2.1, equations (2.1)--(2.3)]{GNII}.
\end{proof}

\subsection{Compatibility with View I and the cross-level identities}
\label{ssec:compat-view-i}

\begin{theorem}[Borcherds--Gritsenko--Nikulin three-way agreement]
\label{thm:bgn-three-way}
The three identifications
\begin{align*}
&\Theta(\phi^{\mathrm{Weil}}_{0,1})=\Delta_5
&&\text{\textup{(}Borcherds 1995, 1998\textup{)}},\\
&64^{-1}\Delta_5(2Z)=\operatorname{den}(\autcor)
&&\text{\textup{(}Gritsenko--Nikulin 1998\textup{)}},\\
&\mathcal D_X=\Delta_5
&&\text{\textup{(}Igusa 1962\textup{)}}
\end{align*}
are mutually consistent:
\(\Theta(\phi^{\mathrm{Weil}}_{0,1})=\mathcal D_X
=64\cdot\operatorname{den}(\autcor)\!\bigm|_{Z\mapsto Z/2}\) on
\(\Hpl^{\IV}_+\), and the multiplier system
\(\nu_{\Delta_5}\) of Theorem~\ref{thm:multiplier-system} is common.
\end{theorem}
\hfill\textnormal{[proved; \cite{Igusa1962,Borcherds95,B,GN,GNII}]}

\begin{proof}
Theorem~\ref{thm:delta5-as-theta-lift} is identification (1).
Theorem~\ref{thm:lift-denominator} is identification (2).
Identification (3) is the Igusa identification
\(\mathcal D_X=\Delta_5\) of \cite[Theorem of \S2]{Igusa1962},
transported through the exterior-square model section.  All three sides have the
Borcherds-product expansion
\(\Delta_5=64\,q^{1/2}r^{1/2}s^{1/2}\prod(1-q^nr^ls^m)^{f(nm,l)}\)
in the cusp chamber and pull back along the exceptional isomorphism of
Theorem~\ref{thm:exceptional-iso} to the same automorphic form on
\(\Sp_4(\Z)\).  The multiplier-system agreement is
Theorem~\ref{thm:multiplier-system}.
\end{proof}

\subsection{The weight-five universal Borcherds identity at \(N=1\)}
\label{ssec:universal-weight-identity}

\begin{theorem}[Universal Borcherds weight identity, \(N=1\)]
\label{thm:universal-weight-identity}
At paramodular level \(N=1\) the Vol III universal Borcherds-weight
identity
\(\kappa_{\mathrm{BKM}}(\Phi_N)=c_N(0)/2\)
specializes to
\[
\kappa_{\mathrm{BKM}}(\Delta_5)
=\frac{c_1(0)}{2}
=\frac{f(0,0)}{2}
=\frac{10}{2}=5.
\]
This is the row \((N=1,\,\operatorname{wt}=5,\,c_1(0)=10)\) of the
Cl\'ery--Gritsenko catalogue \cite{GC}, and is the value entered in
the K3\(\times E\) \(\kappa_\bullet\)-spectrum
\(\{\kappa_{\mathrm{cat}}=0,\kappa_{\mathrm{ch}}^{\mathrm{Hodge}}=0,
\kappa_{\mathrm{ch}}^{\mathrm{Heis}}=3,\kappa_{\mathrm{BKM}}=5,
\kappa_{\mathrm{fiber}}=24\}\) of Vol III.
\end{theorem}
\hfill\textnormal{[proved; \cite{Borcherds95,GC}; Vol III universal
identity]}

\begin{proof}
The lift's weight is \(c_1(0)/2\) by
Theorem~\ref{thm:singular-theta-lift}(1), and
\(c_1(0)=f(0,0)=10\) by Lemma~\ref{lem:phi01-presentations} and
Lemma~\ref{lem:theta-decomposition}.  The
universal identity
\(\kappa_{\mathrm{BKM}}(\Phi_N)=c_N(0)/2\) is
\cite[\S\S2--3]{Gritsenko99}, generalized in Vol III across the eight
Cl\'ery--Gritsenko rows; at \(N=1\) it is exactly the Borcherds
weight formula of Theorem~\ref{thm:singular-theta-lift}(1).
\end{proof}

\begin{remark}[Subscript discipline]
The bare symbol \(\kappa\) does not appear in any theorem-level identity.
Every occurrence carries the subscript fixing its
construction reading
(\(\kappa_{\mathrm{cat}},\,\kappa_{\mathrm{ch}}^{\mathrm{Hodge}},\,
\kappa_{\mathrm{ch}}^{\mathrm{Heis}},\,\kappa_{\mathrm{BKM}},\,
\kappa_{\mathrm{fiber}}\)).  The additive expression
\(\kappa_{\mathrm{BKM}}=\kappa_{\mathrm{ch}}+\chi(\mathcal O_{K3})\) is a
mixed-reading collision.  With
\(\kappa_{\mathrm{ch}}=\kappa_{\mathrm{ch}}^{\mathrm{Heis}}\) it gives the
accidental equality \(3+2=5\) at \(N=1\), but it has no rowwise
extension to \(N\in\{2,3,4,6\}\).  The universal Borcherds-weight identity
\(\kappa_{\mathrm{BKM}}(\Phi_N)=c_N(0)/2\) is the invariant statement;
at \(N=1\) the right hand side is \(c_1(0)/2=5\).
\end{remark}

\subsection{Beilinson cut: the lift is the structural content}
\label{ssec:beilinson-cut-view-v}

The five views of \S\ref{ch:pentadic} lie on the Beilinson chain
\(\mathsf P\rightarrow\mathsf C\rightarrow\mathsf S\rightarrow
\mathsf Z\rightarrow\mathsf A\) of Vol II.  View V's primitive content is
the singular-theta integral
\(\Theta:J_{0,1}^{\mathrm{weak},\mathrm{Weil}}
\rightarrow M_{*}(\widetilde\Orth(\La^{3,2})_+)\) of
Theorem~\ref{thm:singular-theta-lift}, sending Jacobi forms to
Borcherds products.  Its value at \(\phi_{0,1}\) is the automorphic
section \(\Delta_5\), the level-\(\mathsf S\) image of the lift.

\begin{remark}[Beilinson level]
The integral operator \(\Theta\) lives at level \(\mathsf C\) of the
universal chain: it is a chart-dependent passage from a primitive
Jacobi-form input to its Borcherds product, dependent on the
Weil representation of \(\La^{3,2}\) and on the regularization of
\cite[\S6]{B}.  The scalar \(\Delta_5\) lives at level \(\mathsf S\):
it is the automorphic value of the lift on \(\phi_{0,1}\).
The retained \(\mathfrak D_X^{\mathrm{DI}}\) datum of View III is the
conditional Pfaffian construction whose protected Pfaffian recovers
\(\Delta_5\); the BKM superalgebra \(\autcor\) of View II lives at
level \(\mathsf S\) under the trace lane of Vol II.  The primitive
object in this passage is the integral
operator \(\Theta\), and the automorphic, BKM, Pfaffian, and theta-lift
names of \(\Delta_5\) are readings, not primitive objects.
\end{remark}

\subsection{Consistency of the lattice conventions}
\label{ssec:consistency-conventions}

\begin{remark}[\(\La^{3,2}\) versus \(\La^{2,1}\oplus\langle-2\rangle\)]
The signature-\((3,2)\) lattice \(\La^{3,2}\) of
Definition~\ref{def:lambda-3-2} decomposes as
\(\HypPlan\oplus\HypPlan\oplus[2]\) on the integral level.  The
sublattice \(\La^{2,1}=\HypPlan\oplus[2]\) is primitive in
\(\La^{3,2}\); its orthogonal complement is the second hyperbolic
plane \(\HypPlan\), not the rank-one form \(\langle-2\rangle\) (which
would have wrong signature).  The polarization change
\(\alpha:\Gamma_{\mathrm{ind}}\rightarrow\La^{2,1}_{II}\) sends
\((n,l,m)\mapsto 2nf_2-lf_3+2mf_{-2}\), so the BKM grading lattice is
\(\La^{2,1}_{II}\), the index-four sublattice of \(\La^{2,1}\) with
primitive \(f_3\)-coordinate.  The doubled argument \(Z\mapsto 2Z\) in
\(\operatorname{den}(\autcor)=64^{-1}\Delta_5(2Z)\) of
Gritsenko--Nikulin changes the exponential convention from
\(\exp(-\pi i(\alpha,z))\) to \(\exp(-2\pi i(\alpha,z))\).  The
\(\mathrm{II}_{2,2}\) presentation of \(\La^{3,2}\) used in some
references (Borcherds 1995 \cite[\S15]{Borcherds95},
Borcherds 1998 \cite[Example 16.4]{B}) is the same lattice expressed via two
hyperbolic planes plus the \([2]\)-form, and produces the same Borcherds
product \(\Delta_5\) under the same lift.  No
\(\mathrm{II}_{3,2}\) versus \(\mathrm{II}_{2,1}\oplus\langle-2\rangle\)
ambiguity remains.
\end{remark}

\begin{remark}[Multiplier and metaplectic accounting]
The Borcherds singular-theta lift produces an automorphic form on the
integral form \(\widetilde\Orth(\La^{3,2})_+\), pulled back to
\(\Sp_4(\Z)\) by Theorem~\ref{thm:exceptional-iso}.  Since
\(c_1(0)/2=5\) is integral, the resulting automorphy factor is defined on
\(\Sp_4(\Z)\).  The metaplectic cover enters
on the \(\Mp_2(\Z)\) input side through
\(\rho_{\La^{3,2}}\); after restriction to integral weight its
central sign does not alter the \(\Sp_4(\Z)\) automorphy factor.  The
residual reflection character is exactly the
multiplier \(\nu_{\Delta_5}\) of Theorem~\ref{thm:multiplier-system}.
For the CHL constants at \(N\in\{1,2,3,6\}\), the numbers
\(c_N(0)/2=(5,3,2,1)\) are integral, and the metaplectic cover
belongs to the input Weil representation while the resulting
automorphic forms live on the corresponding integral orthogonal or
Siegel arithmetic groups.  At \(N=4\) the constant \(c_4(0)=3\) is
odd and the weight \(c_4(0)/2=\tfrac32\) is half-integral
(Govindarajan--Krishna \cite{GK10Composite}); \(\Phi_4\) is a genuine
form on the genus-two metaplectic cover.
\end{remark}

\begin{remark}[The eta-multiplier \texorpdfstring{\(v_\eta^{12}\)}{v\_eta\^12}]
The character of \(\Delta_5\) on the Jacobi parabolic of \(\Sp_4(\Z)\)
is described in \cite{Gritsenko99} as \(v_\eta^{12}\times v_H\), where
\(v_\eta\) is the eta-multiplier on \(\SL_2(\Z)\) and \(v_H\) is the
Heisenberg multiplier.  The first four Cl\'ery--Gritsenko rows carry
eta exponents \(12,6,4,3\) for \(N=1,2,3,4\), respectively
\cite[Theorem 1.4]{GC}; the present form is the \(N=1\) row.
Equivalently, after the lift through
the exceptional isomorphism, \(v_\eta^{12}\) remains the eta multiplier
in the Jacobi--Heisenberg presentation; it is not a discriminant
character of \(\mathrm{II}_{2,2}\), since \(\mathrm{II}_{2,2}\) is
unimodular.  No \(v_\eta^N\) inconsistency between
\cite{Gritsenko99}, \cite{GN}, and \cite{GC} arises at \(N=1\); the
three computations agree.
\end{remark}

\subsection{Theta-lift boundary}
\label{ssec:theta-lift-boundary}

The singular theta lift is one realization of \(\Delta_5\) among five.
The lift's primitive content is the Borcherds regularized theta
operator \(\Theta:\phi^{\mathrm{Weil}}\mapsto\Theta(\phi^{\mathrm{Weil}})\)
of Theorem~\ref{thm:singular-theta-lift}; its value at
\(\phi^{\mathrm{Weil}}_{0,1}\) is the weight-five Igusa cusp form.  The five-view
identification
\(\mathcal D_X=64\cdot\operatorname{den}(\autcor)\!\bigm|_{Z\mapsto Z/2}
=\Theta(\phi^{\mathrm{Weil}}_{0,1})\) is the theorem-level
automorphic--BKM--theta identity; its Pfaffian extension to View III
is conditional on the retained \(K3\times E\) data of
Appendix~\ref{app:obstruction-discharges}.  The mirror-discriminant View IV is the
remaining conjectural cross-level identity, recorded in the
mirror-discriminant chapter (Chapter~8) and not used here.

Borcherds's theta-correspondence formulation makes two structural facts visible.
First, the weight-five integrality is forced by
\(c_1(0)=10\) and the integrality of \(c_1(0)/2\); no metaplectic cover
is needed at the weight level.  Second, the BKM-denominator structure
of View II is not an extra postulate but a theorem about the lift's
image: the Gritsenko--Nikulin denominator identity reads the lift's
Borcherds product as the denominator of \(\autcor\) on
\(\La^{2,1}_{II}\) (Theorem~\ref{thm:lift-denominator}), and the
Borcherds--Gritsenko--Nikulin three-way agreement
(Theorem~\ref{thm:bgn-three-way}) crystallizes the lift's value in
three convertible normalizations.  The universal identity
\(\kappa_{\mathrm{BKM}}(\Phi_N)=c_N(0)/2\) is, at
\(N=1\), precisely the Borcherds-weight formula
\(\operatorname{wt}\Theta(\phi)=c_\phi(0,0)/2\) of
Theorem~\ref{thm:singular-theta-lift}(1) applied to
\(\phi=\phi_{0,1}\).
